\section{三维贯通实验：从 Token ID 到三个位置的输出概率}
\label{intro:22b-token-to-output}

之前已经分别认识 Embedding、Q/K/V、因果 Mask、位置、Pre-Norm、两条残差、FFN 与输出头。
但这些组件若一直用不同的玩具数字来算，读者就无法判断它们连起来是否真的一致。
现在\textbf{只选一套三维参数}，把一层 Decoder-only Transformer
从输入走到每个位置的五类概率。\textbf{以下全部数值来自仓库
\texttt{scripts/teaching\_reference.py}，而不是临时另外设定的 logits。}

\subsection{第一步：Token ID、位置向量与形状}

词表按 ID 0--4 依次是“其他、春天、。、花、开了”。
整句的 ID 是 $[1,3,4,2]$；本次输入取前三个 $[1,3,4]$，
正确的三位置目标是 $[3,4,2]$。Embedding 使用前面唯一的 $5\times3$ 初值表：
\[
E_{[1,3,4]}=
\begin{bmatrix}
1.2&0&-0.2\\-0.3&1.1&0.2\\0.15&-0.2&1.25
\end{bmatrix}.
\]
为区分同一词在不同位置，\textbf{固定的位置向量}是
\[
P=\begin{bmatrix}
0.10&0.05&-0.10\\0&-0.05&0.12\\-0.12&0.10&0.07
\end{bmatrix}.
\]
加到各自对应行后，Block 的输入为
\begin{equation}
\boxed{X=E_{[1,3,4]}+P=
\begin{bmatrix}
1.30&0.05&-0.30\\-0.30&1.05&0.32\\0.03&-0.10&1.32
\end{bmatrix}\in\R^{3\times3}.}
\label{eq:running-three-token-x}
\end{equation}
位置是输入的一部分，不等于 Attention 的 Mask。
两行即使有相同 Token ID，也会因为对应位置向量不同而得到不同的 $X$。

\subsection{第二步：先 Norm，再从三个位置算 Q/K/V}

用\textbf{每个 Token 行独立}的 LayerNorm：
\[
U=\operatorname{LN}(X),\qquad
\operatorname{LN}(\vect x^\top)=
\frac{\vect x^\top-\mu}
{\sqrt{\frac{1}{3}\sum_j(x_j-\mu)^2+10^{-5}}}.
\]
这个教学网络把 LayerNorm 的缩放 $\gamma=1$、平移 $\beta=0$ \textbf{固定}，
投影和 FFN 都无偏置；它们不是一般 Transformer 必须采用的限制。
得到
\[
U\approx\begin{bmatrix}
1.3833&-0.4368&-0.9464\\
-1.1901&1.2566&-0.0665\\
-0.6033&-0.8061&1.4093
\end{bmatrix}.
\]
独立的可训练投影 $Q=UW_Q,\ K=UW_K,\ V=UW_V$：
\[
W_Q=\begin{bmatrix}.8&.2&-.1\\.1&.9&.25\\-.1&.15&.8\end{bmatrix},\quad
W_K=\begin{bmatrix}.85&0&.2\\-.1&.9&0\\.15&.2&.75\end{bmatrix},\quad
W_V=\begin{bmatrix}1&.2&0\\0&.9&.15\\.1&-.2&.85\end{bmatrix}.
\]
例如三行投影结果为
\[
Q\approx\begin{bmatrix}
1.1576&-0.2584&-1.0047\\
-0.8198&0.8829&0.3800\\
-0.7041&-0.6347&0.9863
\end{bmatrix},\qquad
K\approx\begin{bmatrix}
1.0775&-0.5824&-0.4332\\
-1.1473&1.1176&-0.2879\\
-0.2208&-0.4436&0.9364
\end{bmatrix}.
\]
\[
V\approx\begin{bmatrix}
1.2886&0.0728&-0.8700\\
-1.1968&0.9062&0.1320\\
-0.4623&-1.1280&1.0770
\end{bmatrix}.
\]
这一次 $W_Q,W_K,W_V$ 确实不同，并不采用 $Q=K=V=X$ 的退化捷径。

\subsection{第三步：因果可见性、Softmax 与 Value 读取}

三个位置有 $d_k=3$，所以分数除以 $\sqrt3$。
再将不可见的未来位置设为 $-\infty$：
\begin{equation}
S=\frac{QK^\top}{\sqrt3}+C
\approx
\begin{bmatrix}
1.0583&-\infty&-\infty\\
-0.9019&1.0496&-\infty\\
-0.4713&-0.1071&0.7855
\end{bmatrix}.
\label{eq:two-token-scores}
\end{equation}
逐行 Softmax 得到读取比例：
\begin{equation}
\boxed{A\approx\begin{bmatrix}
1&0&0\\
0.1244&0.8756&0\\
0.1680&0.2418&0.5903
\end{bmatrix}.}
\label{eq:three-token-attention}
\end{equation}
例如第三行只在前三个已知位置中分配权重；第二行不能读取“开了”。
按权重读出 Value，$R_{\rm raw}=AV$：
\[
R_{\rm raw}\approx\begin{bmatrix}
1.2886&0.0728&-0.8700\\
-0.8876&0.8025&0.0074\\
-0.3458&-0.4345&0.5215
\end{bmatrix}.
\]
随后用\textbf{注意力输出投影} $W_O$，再走第一条残差：
\[
W_O=\begin{bmatrix}.9&0&.1\\.1&.85&0\\0&.15&.95\end{bmatrix},
\qquad
X_{\rm att}=X+R_{\rm raw}W_O
\approx\begin{bmatrix}
2.4670&-0.0186&-0.9976\\
-1.0186&1.7333&0.2382\\
-0.3247&-0.3911&1.7809
\end{bmatrix}.
\]
注意残差相加的两边都是 $3\times3$，并没有把 Token 行数改变。

\subsection{第四步：第二次 Norm、升维 FFN 与残差}

让 $B=\operatorname{LN}(X_{\rm att})$，再计算
\[
F_{\rm mid}=\relu(BW_{\rm up}),\qquad
H=X_{\rm att}+F_{\rm mid}W_{\rm down},
\]
其中 $W_{\rm up}:3\times4$，$W_{\rm down}:4\times3$。
两套矩阵的初值如下（其余维度不能随意替换）：
\[
W_{\rm up}=\begin{bmatrix}
.4&.2&-.3&.1\\-.1&.5&.15&-.25\\.2&.1&.45&.3
\end{bmatrix},\quad
W_{\rm down}=\begin{bmatrix}
.25&-.1&.15\\-.05&.2&.1\\.1&.15&.15\\.1&-.12&.18
\end{bmatrix}.
\]
归一化和 FFN 激活结果可分别核对：
\[
B\approx\begin{bmatrix}
1.3601&-0.3444&-1.0157\\
-1.1879&1.2585&-0.0706\\
-0.6739&-0.7398&1.4137
\end{bmatrix},
\quad
F_{\rm mid}\approx
\begin{bmatrix}
0.3753&0&0&0\\
0&0.3846&0.5134&0\\
0.0872&0&0.7274&0.5417
\end{bmatrix}.
\]
得到 Block 最终三行隐藏状态：
\begin{equation}
\boxed{H\approx\begin{bmatrix}
2.5609&-0.0561&-0.9413\\
-0.9865&1.8872&0.3537\\
-0.1760&-0.3557&2.0005
\end{bmatrix}\in\R^{3\times3}.}
\label{eq:three-token-block-output}
\end{equation}

\subsection{第五步：共享的输出头一次生成三个位置的概率}

用同一个词表头矩阵
\[
W_{\rm out}=\begin{bmatrix}
0&.2&0&.2&-.2\\0&.2&0&1.2&-.4\\0&-.1&.4&.1&.2
\end{bmatrix}\quad (3\times5)
\]
计算 $Z=HW_{\rm out}$，并逐行 Softmax：
\[
Z\approx
\begin{bmatrix}
0&0.5951&-0.3765&0.3507&-0.6780\\
0&0.1448&0.1415&2.1027&-0.4868\\
0&-0.3064&0.8002&-0.2620&0.5776
\end{bmatrix}.
\]
按 ID 0--4 排列的三行概率是
\begin{equation}
\boxed{\widehat P\approx\begin{bmatrix}
0.1843&0.3341&0.1264&0.2617&0.0935\\
0.0826&0.0954&0.0951&0.6761&0.0507\\
0.1535&0.1130&0.3418&0.1181&0.2735
\end{bmatrix}\in\R^{3\times5}.}
\label{eq:three-token-probabilities}
\end{equation}
第一行最大候选是 ID 1“春天”，目标却是 ID 3“花”；
第二行最大候选是 ID 3“花”，目标却是 ID 4“开了”；
第三行最大候选是 ID 2“。”，恰好等于目标。
\textbf{训练前两题错、一题对}，并非人工指定的每题都错。

\begin{center}\small
\tagpdfsetup{table/header-rows={1}}
\begin{tabularx}{\linewidth}{@{}p{29mm}p{27mm}p{26mm}X@{}}
\toprule
阶段 & 输入形状 & 输出形状 & 新增的主要能力\\\midrule
Embedding + 位置 & $3\times3$ & $3\times3$ & 身份与位置\\
Norm + Q/K/V & $3\times3$ & 各 $3\times3$ & 独立投影\\
Mask + Softmax & $3\times3$ & $3\times3$ & 不读取未来\\
Value + 输出投影 & $3\times3$ & $3\times3$ & 跨位置读取\\
第二次 Norm + FFN & $3\times3$ & $3\times3$ & 逐位置非线性\\
词表输出头 & $3\times3$ & $3\times5$ & 每位置五类分数\\
\bottomrule
\end{tabularx}
\end{center}

\subsection{换一项输入，最后一行会怎样？}

把第一个输入由 ID 1 换为 ID 4，得到 $[4,3,4]$，
保持最后两个 Token 和所有模型参数不变，
第三行五个 logits 从
$(0,-0.3064,0.8002,-0.2620,0.5776)$
变为约 $(0,-0.4279,0.9325,-0.6140,0.7915)$。
最大分数差约 $0.3520$。
这是\textbf{最后位置利用了前文}的可复算对照；它不是训练或重新拟合。

\begin{warningbox}
这是完整的\textbf{一层单头、无偏置、固定 Norm 仿射参数}的教学 Block，
不等于包含多层、多头及所有工程配置的真实 GPT。
其他章节的二维/四维局部例子已经单独标明；本章只使用三维贯通参数。
\end{warningbox}

\par\medskip\noindent{\sffamily\footnotesize\color{blue}\textbf{读完自测}}\quad
\phantomsection\label{check:22b-token-to-output}
统一三维例的三个位置中，哪两个初始贪心预测错误？如果将第一 Token 替换为 ID 4，第三位置 logits 是否完全不变？
\quad{\footnotesize\hyperref[answer:22b-token-to-output]{核对解答。}}\par\medskip

\begin{conclusion}
从 ID 到三位置概率的完整链已打通：每一个中间量都能通过同一参数表复算。
下一个缺口：这三份预测怎样\textbf{同时}给所有共享参数提供训练梯度？
\end{conclusion}
